% Français 9115, batch 3 — Gallica views 41–60.
% Pass: Opus 5 (claude-opus-5), 2026-09-12 — first pass, unchecked against
% the views by a human.
%
% What the twenty views hold. Batch 3 closes one piece and opens another.
% View 41 carries the last eight lines of the « Exposition des principes qui
% servent de base à la théorie des surfaces élastiques », which batches 1 and
% 2 transcribed from view 3 onward: the sentence view 40 broke off in
% finishes, a horizontal rule is drawn under it, and the piece ends. View 42
% is that leaf's blank verso — the recto's writing shows through, reversed,
% and nothing else — and is not transcribed.
%
% Everything from view 43 to view 60 is a second fair copy, in the same
% regular hand, with its own pagination running 1 to 16 and no break: five
% pages of « Observations préliminaires » (v. 43–47) followed by the
% « Mémoire sur la courbure des surfaces » proper (v. 48–60), which is still
% running at the foot of view 60. This is the memoir on mean curvature, and
% the subject is no longer elasticity but the geometry underneath it: a
% quantity she says has not been noticed, the mean curvature of a surface at
% a point, and the law by which curvature is distributed around that point.
%
%   v. 41      End of the Exposition. « Si par un hazard singulier j'avois
%              rencontré juste, il seroit fâcheux qu'une semblable défaveur
%              détournât leur attention d'une théorie, si simple, si
%              générale, si naturelle ». Rule; end of piece.
%   v. 42      Blank verso. Not transcribed.
%   v. 43–47   « Observations préliminaires », her pages 1–5. What she fears
%              is not the dryness of the subject but the novelty of the
%              angle: she means to name and define « un genre de quantités
%              dont l'existence ne paroît pas avoir été soupçonné ». Then the
%              case for mean curvature — a force born of curvature may depend
%              on the mean value and not on the distribution, so two surfaces
%              of wholly different aspect can carry equal quantities of
%              curvature; the analogy with the distribution of heat in the
%              ring is announced (v. 45) and defended as not fortuitous,
%              both resting on a circular space; and the construction is
%              stated (v. 46–47) that makes curvature proportional to the
%              convex surface of a small cylinder, elliptic-based where the
%              principal radii lie the same way. Page 5 breaks off in
%              mid-sentence at the foot.
%   v. 48–49   N° 1. The memoir opens on Euler — « Recherches sur la courbure
%              des surfaces » (M. de Berlin 1760 p. 119), quoted at length,
%              including the objection she has to answer: the curvature of a
%              surface is « fort équivoque », there being an infinity of
%              curvatures at each point, and Euler's own remedy, the normal
%              sections alone. Then her own history of the question: the
%              objection was made to her when she proposed the sum of the
%              inverse ratios of the two principal radii; what she did in the
%              memoir sent to the Académie « pour le concours de 1816 »; and
%              that the proof « fut formellement désapprouvée » (v. 50).
%   v. 50–51   The invariance she now rests on: for any two perpendicular
%              normal planes the sum of the inverse radii equals that of the
%              principal ones; the four quadrants of the sphere of that
%              radius lie alternately outside and inside the surface; and the
%              compensation between curvatures greater and smaller than the
%              sphere's. Then the decision to treat the question on its own,
%              apart from any application, and N° 2 begins.
%   v. 52      A slip (béquet) photographed on its own: see the note at that
%              view.
%   v. 53      A second photograph of the page of view 51. Not transcribed
%              twice.
%   v. 54      N° 2 carried out. Euler's formula for the radius of an oblique
%              normal section and its perpendicular, and the two curvatures
%              in the form 1/r = (f sin²φ + g cos²φ)/fg,
%              1/r′ = (f cos²φ + g sin²φ)/fg. Then the objection she raises
%              against herself: the subject is the curvature of surfaces and
%              the formulas employed are those of linear curvature.
%   v. 55–60   N° 3, the law of the distribution of curvature, set beside
%              Fourier's law for the ring, proposition by proposition. The
%              one real difference is announced first (v. 55): a quarter
%              circumference separates the extreme curvatures where a half
%              circumference separates the extreme temperatures. Then
%              1st prop. (v. 56), the sum in two perpendicular normal planes
%              is invariant, with Fourier cited (Théorie de la chaleur,
%              N° 247, p. 277–279); 2nd prop. (v. 57–58), the section at 45°
%              carries the mean curvature, with the long quotation from
%              Fourier on the two points of the ring at the mean temperature;
%              her own restatement « par rapport à la courbure » (v. 59); and
%              3rd prop. (v. 60), the four directions around the point along
%              which the normal-plane curves are alike, reduced to two for
%              the principal curvatures. The batch ends mid-sentence.
%
% The numbering on the leaves, and why no \folio{} is written. The ink series
% of batches 1 and 2 reaches 21 on view 41, with the struck second number the
% earlier batches recorded. From view 43 the series continues 22, 23, 24, 25,
% 26, 27, 28, 29 on the rectos — views 43, 45, 47, 49, 51 (and 53, which is
% the same leaf again), 55, 57, 59 — and is joined by a second number which
% is the author's own pagination: odd beside the ink number on the recto
% (1, 3, 5, 7, 9, 11, 13, 15), even in the top left corner of the verso
% (2, 4, 6, 8, 10, 12, 14, 16). Every one of those numbers was read on the
% view except the « 1 » of view 43, which is barely formed. None is inferred.
% They are in ink and not in the Bibliothèque's pencil; no pencilled
% foliation was found anywhere in the batch, so, as in batches 1 and 2, no
% \folio{} is written.
%
% Notation. f and g are the greatest and least radii of principal curvature;
% r and r′ the osculating radii of the curves in two perpendicular normal
% planes; φ the angle the normal plane makes with that of greatest curvature,
% φ′ the auxiliary angle of N° 3. She writes the trigonometric functions with
% a full stop — « cos.φ′ », « sin.2φ′ » — and that stop is kept here. Her
% « demi somme » is written in two words. Two places where the leaf does not
% say what the sense wants are transcribed as she wrote them and noted at the
% spot: the quantity called the half-sum on view 57, written 1/f + 1/g, and
% the four identities of view 60, written with the factor 1/2 where the
% calculation that follows them uses √2/2. Neither was corrected.
%
% What could not be read. Seven places: the abbreviated numeral after
% « Enfin » on view 51; an ink blot swallowing one word on view 50; the
% struck and overwritten cluster after « mêlées » and the second one at the
% end of that sentence on view 51; the one or two struck words before « plan
% normal » on view 54; the abbreviation before « IV »
% in the Fourier reference on view 56, whose line also runs off the edge of
% the leaf, taking the last digits of the page range with it; and the struck
% words after « dans les lois que nous » on view 59. None of these was
% guessed. Her own underlining is set \emph{}; \uncertain{} renders as an
% underline too, so a reader of the PDF should take this header's word for
% which is which. In the quotations from Euler and from Fourier she repeats a
% guillemet at the head of every line; as in batch 2 only the opening and
% closing ones are kept here.
%
% Nothing here was checked against any published transcription.
\documentclass[11pt,a4paper]{article}
% The preamble shared by every transcription of the mathematicians' manuscripts in Gallica.
%
% It rests on one idea: **uncertainty compounds, it does not resolve.** A
% transcription that smooths over a doubtful word has destroyed the only thing
% it was made for — a reader must be able to tell, at every point, what was on
% the page and what was supplied. The apparatus macros below are therefore the
% critical apparatus, not decoration.
%
% The same commands are recognised by `scripts/render.mjs`, which produces the
% site's reading view, and by `scripts/tei.mjs`. A macro added here must be
% added there too, or rendering fails loudly — which is the wanted behaviour.
%
% Comments here are English, like the rest of the repository. The documents
% this preamble sets are French, and that is deliberate: see the skills.

\ProvidesPackage{germain}[2026/09/15 Transcriptions of mathematicians' manuscripts in Gallica]

\usepackage{amsmath,amssymb,amsthm}
\usepackage{mathrsfs}
% Kept from the parent preamble so the renderer's diagram support stays
% available; these manuscripts draw no commutative diagrams, and nothing here uses it.
\usepackage{tikz-cd}
\usepackage[english,french]{babel}
\usepackage{fontspec}

% --- Greek ------------------------------------------------------------------
% The text face stays fontspec's default, Latin Modern, which has no Greek:
% XeTeX drops every glyph it lacks with a line in the log and nothing on the
% page, and the Greek words the manuscripts quote vanished from the PDFs. So
% the Greek and Greek Extended blocks get a XeTeX character class of their own,
% and entering or leaving a run of it switches the family to CMU Serif — the
% same Computer Modern design, with polytonic Greek — and back. Only the
% family changes: a Greek word in \textit or \textbf keeps its shape and series.
% The transitions to and from every other class carry the tokens that class
% has towards an ordinary letter, so French spacing before « ; » or inside
% « » still holds after a Greek word. No grouping, so no brace or \endgroup
% can come out unbalanced; maths is left alone, since XeTeX inserts nothing
% there. Kept identical in `exercices.sty`.
\makeatletter
\newfontfamily\germainGreekfont{cmun}[
  Extension=.otf, NFSSFamily=germaingreek,
  UprightFont=*rm, ItalicFont=*ti, SlantedFont=*sl,
  BoldFont=*bx, BoldItalicFont=*bi, BoldSlantedFont=*bl]
\newXeTeXintercharclass\germain@greek
\def\germain@greekrange#1#2{%
  \@tempcnta=#1\relax
  \loop
    \XeTeXcharclass\@tempcnta=\germain@greek
    \ifnum\@tempcnta<#2\advance\@tempcnta\@ne\repeat}
\germain@greekrange{"0370}{"03FF}
\germain@greekrange{"1F00}{"1FFF}
\def\germain@greekprev{\rmdefault}
\def\germain@greekin{%
  \edef\germain@greekprev{\f@family}\fontfamily{germaingreek}\selectfont}
\def\germain@greekout{\fontfamily{\germain@greekprev}\selectfont}
\def\germain@greekjoin{%
  \XeTeXinterchartokenstate=\@ne
  \@tempcnta=\z@
  \loop
    \ifnum\@tempcnta=\germain@greek\else
      \XeTeXinterchartoks\germain@greek\@tempcnta=\expandafter{%
        \expandafter\germain@greekout\the\XeTeXinterchartoks\z@\@tempcnta}%
      \XeTeXinterchartoks\@tempcnta\germain@greek=\expandafter{%
        \the\XeTeXinterchartoks\@tempcnta\z@\germain@greekin}%
    \fi
    \ifnum\@tempcnta<4095\advance\@tempcnta\@ne\repeat}
% babel-french sets its own classes, and saves and restores the interchar
% state, each time a language is selected: join again after it does.
\addto\extrasfrench{\germain@greekjoin}
\addto\extrasenglish{\germain@greekjoin}
\AtBeginDocument{\germain@greekjoin}
\makeatother

\usepackage{geometry}
\geometry{a4paper,margin=25mm,marginparwidth=28mm}
\usepackage{marginnote}
\usepackage{xcolor}
\usepackage[normalem]{ulem}
\setlength{\emergencystretch}{3em}
\usepackage{hyperref}
\hypersetup{colorlinks=true,linkcolor=black,urlcolor=black,pdfborder={0 0 0}}

% --- Batch metadata ---------------------------------------------------------
% Taken as they stand by the rendering script, which shows them at the head of
% the reading view. They fix what the file claims to cover. The macro names are
% the parent project's, so that the scripts stay shared; \folder here means the
% volume's slug — `fr-9115` — and \pages counts Gallica views.

\newcommand{\folder}[1]{\gdef\grothFolder{#1}}
\newcommand{\batch}[1]{\gdef\grothBatch{#1}}
\newcommand{\pages}[2]{\gdef\grothFirst{#1}\gdef\grothLast{#2}}
\newcommand{\foldertitle}[1]{\gdef\grothFoldertitle{#1}}

% The holder's own shelfmark, printed as they write it: « Français 9115 ».
\newcommand{\shelfmark}[1]{\gdef\grothShelfmark{#1}}

% The dating, copied verbatim from the holder's catalogue — never inferred
% here. « XIXe siècle » is the BnF's; a pass that dates a leaf from its content
% says so in a \note{}, as its own inference.
\newcommand{\dating}[1]{\gdef\grothDating{#1}}

% The status notice, printed in the title block and repeated as a diagonal
% watermark on every page of the PDF. The manuscripts are in the public
% domain; what this line declares is not a right but a quality — an
% unverified first pass by a machine. The skills write the line; a file
% without it gets no watermark, so leaving it out is visible in the output.
\newcommand{\watermark}[1]{\gdef\grothWatermark{#1}}

\gdef\grothFolder{?}\gdef\grothBatch{}\gdef\grothFirst{?}\gdef\grothLast{?}%
\gdef\grothFoldertitle{}\gdef\grothShelfmark{}\gdef\grothDating{}\gdef\grothWatermark{}

% --- The résumé -------------------------------------------------------------
\newenvironment{resume}
  {\small\setlength{\parskip}{0.45em}}
  {\par\medskip\hrule\medskip}

% --- Keywords ---------------------------------------------------------------
% In English, closing the modernised reading's résumé: the modern vocabulary
% under which to search for what the leaves construct. The manifest extracts
% them to tag the volume — this line is the single source of the tags.
\newcommand{\keywords}[1]{\par\smallskip\noindent
  {\footnotesize\textsc{Keywords} --- \textit{#1}}\par}

% --- The critical apparatus -------------------------------------------------

% The Gallica view number, in the margin. This is the anchor that lets a
% disagreement over a formula be settled by looking at *one* image rather than
% a volume of 750. The site uses it to turn the facsimile as the transcription
% is scrolled. A view is one scanned image: usually one side of a leaf.
\newcommand{\page}[1]{%
  \marginnote{\footnotesize\color{gray!70}\textsf{v.\,#1}}%
  \label{page:#1}\ignorespaces}

% The BnF's pencilled foliation, where the leaf carries it and a pass can read
% it: « 198r ». This is what the literature cites — « ff. 198r–208v » — and
% the one line that turns a citation into a view. Recorded, never inferred:
% a folio a pass cannot read is not written.
\newcommand{\folio}[1]{%
  \marginnote{\footnotesize\color{gray!50}\textsf{f.\,#1}}\ignorespaces}

% The modernised reading's coarser counterpart to \page{}: which views a
% section is drawn from.
\newcommand{\pagerange}[2]{%
  \marginnote{\footnotesize\color{gray!70}\textsf{v.\,#1--#2}}%
  \ignorespaces}

% Illegible: never guessed.
\newcommand{\ill}{\textcolor{red!60!black}{[\,\ldots\,]}}

% A doubtful reading: the word is given, and flagged as uncertain.
\newcommand{\uncertain}[1]{\uline{#1}}

% An editorial addition — a missing letter, an implied word.
\newcommand{\add}[1]{\textcolor{blue!50!black}{[#1]}}

% Struck out by the writer. What was crossed out is part of the document.
\newcommand{\struck}[1]{\sout{#1}}

% The transcriber's note, on the state of the page or the mathematics.
\newcommand{\note}[1]{\par\smallskip{\small\color{gray!60!black}\textit{[#1]}}\par\smallskip}

% A marginal note of hers, distinct from the body of the text.
\newcommand{\marginal}[1]{\par\smallskip\noindent\hspace{1em}%
  {\small\color{gray!60!black}#1}\par\smallskip}

% --- Title ------------------------------------------------------------------
% The title block is French, like the documents themselves.

\AddToHook{shipout/background}{%
  \ifx\grothWatermark\empty\else
    \begin{tikzpicture}[remember picture, overlay]
      \node[rotate=52, scale=3.4, align=center, text=gray!18,
            font=\sffamily\bfseries]
        at (current page.center) {\grothWatermark};
    \end{tikzpicture}%
  \fi}

\AtBeginDocument{%
  \begin{center}
    {\large\bfseries Manuscrits de mathématiciens (Gallica) ---
      \ifx\grothShelfmark\empty\grothFolder\else\grothShelfmark\fi}\\[2pt]
    {\itshape\grothFoldertitle}\\[4pt]
    {\small vues \grothFirst--\grothLast
      \ifx\grothBatch\empty\else\ (lot \grothBatch)\fi}\\[2pt]
    \ifx\grothDating\empty\else
      {\small\color{gray!60!black}Datation du catalogue : \grothDating}\\[2pt]
    \fi
    \ifx\grothWatermark\empty\else
      {\small\bfseries\color{red!45!gray}\grothWatermark}\\[2pt]
    \fi
    {\footnotesize\color{gray!60!black}Transcription établie d'après les images IIIF de Gallica.
     Source gallica.bnf.fr / Bibliothèque nationale de France.\\
     Les lectures incertaines sont soulignées, les illisibles notées [\,\ldots\,],
     les ajouts entre crochets ; \textsf{v.} numérote les vues de Gallica,
     \textsf{f.} la foliotation au crayon de la BnF.}
  \end{center}
  \vspace{1em}}

\endinput


\folder{fr-9115}
\shelfmark{Français 9115}
\batch{3}
\pages{41}{60}
\dating{1801-1900}
\watermark{Lecture automatique\\non vérifiée}
\foldertitle{Papiers de Sophie GERMAIN. — Recueil de dissertations et problèmes mathématiques et physiques.}

\begin{document}

\note{Numérotation des feuillets. La série à l'encre des lots 1 et 2 atteint
21 à la vue 41, accompagnée du second nombre biffé déjà relevé. À partir de
la vue 43 elle se poursuit sans interruption — vues 43, 45, 47, 49, 51 (et
53, qui est le même feuillet), 55, 57, 59 → 22, 23, 24, 25, 26, 27, 28,
29 — et s'accompagne d'un second nombre, qui est la pagination de l'auteur :
impair à côté du nombre à l'encre au recto (1, 3, 5, 7, 9, 11, 13, 15), pair
en haut à gauche des versos (2, 4, 6, 8, 10, 12, 14, 16). Tous ces nombres
ont été lus sur la vue, sauf le « 1 » de la vue 43, à peine formé. Ils sont à
l'encre et non au crayon : aucune foliotation de la Bibliothèque n'a été lue
sur ces vues, et aucun « f. » n'est donc porté dans cette transcription.}

\page{41}
ne dédaignent d'examiner une opinion qui n'est encore que la mienne.

Si par un hazard singulier j'avois rencontré juste, il seroit fâcheux qu'une
semblable défaveur détournât leur attention d'une théorie, si simple, si
générale, si naturelle, d'une théorie enfin qui se réduiroit à n'être que
l'énoncé même de la question, et à laquelle il n'auroit manqué, pour paroître
susceptible des plus belles applications, que d'être tombée en des mains plus
habiles

\note{un trait horizontal ferme la page : c'est la fin de l'« Exposition des
principes qui servent de base à la théorie des surfaces élastiques »,
commencée à la vue 3. La phrase n'a pas de point final. La vue 42 est le
verso resté blanc de ce feuillet — seule l'écriture du recto y transparaît,
inversée — et n'est pas transcrite.}

\section{Observations préliminaires}

\page{43}
Des surfaces ont été l'objet d'un grand nombre de travaux ; et cependant,
quand je cherche à préjuger quel obstacle pourroit, à la première vue,
décourager ici le lecteur, je suis moins effrayée de l'espèce de sécheresse
qui s'attache au sujet, que de la trop grande nouveauté de l'aspect sous
lequel j'ai crû devoir l'envisager.

Je me propose, en effet, de signaler et de définir un genre de quantités dont
l'existence ne paroît pas avoir été soupçonné ; et, en traitant une matière
qui semble être épuisée, j'aurai\struck{s} recours à des distinctions
inusitées. N'est-il pas à craindre que des notions singulières, en apparence,
soient repoussées avant l'examen, lorsqu'elles n'ont pour appui ni l'autorité
d'un nom célèbre, ni, malgré tous mes soins, le secours d'une exposition
lumineuse.

Une observation préliminaire paroît toutefois propre à faire pressentir, au
moins, que la considération des courbures moyennes, objet principal des
recherches qui vont être exposées, n'est peut être pas tout à fait sans
utilité.

Tant qu'on se borne à examiner les surfaces par rapport à elles seules, la
nature des choses établit une liaison tellement intime entre l'idée de la
figure et celle de la courbure que les mots qui expriment ces idées sont,
dans plusieurs cas, parfaitement synonimes ; mais il n'en est plus de même
lorsque la courbure entre en comparaison avec des quantités dynamiques, comme
il arrive dans certaines questions. Alors la courbure est

\page{44}
tacitement traitée comme une quantité du même genre ; suivant les cas, il
naît de cette quantité, soit une résistance soit une force ; on peut donc
accorder, d'abord, que la courbure soit susceptible d'une valeur moyenne. De
plus, si on s'attache à considérer un point donné de la surface, on trouve, à
la vérité, que la figure dépend de la répartition de la courbure autour de ce
point ; mais il ne répugne en aucune manière d'admettre que la force née de
la courbure dépende, non de la répartition, mais de la valeur moyenne de
cette courbure et que, par conséquent, en ce sens, deux surfaces qui
présentent à l'œil des aspects entièrement différens possèdent cependant,
dans les points comparés, des quantités égales de courbures.

On conçoit peut être déjà comment, dans une question où l'évaluation des
forces dépendroit de la courbure des surfaces, l'existence des courbures
moyennes pourroit offrir de précieuses simplifications ; mais l'aperçu de
cette espèce de convenance ne suffit pas ; il faut encore que l'existence
effective de la quantité qui en est l'objet soit établie d'une manière
directe.

Or, si la considération de la courbure des surfaces peut se mêler à d'autres
sujets, elle ne cesse pourtant pas d'être essentiellement du domaine de la
géométrie. L'existence des courbures moyennes doit donc être établie d'une
manière géométrique.

\note{sur cette page plusieurs mots du dernier paragraphe — « elle ne cesse »,
« essentiellement », « géométrie » — sont repassés d'une encre plus grasse,
sans rature : c'est une reprise du trait et non une correction.}

\page{45}
On a peu ajouté à ce qu'Euler nous a appris touchant la courbure des
surfaces. Je prendrai pour point de départ le point où ce grand géomètre
s'est arrêté : et, en m'aidant de ses travaux, j'essayerai de pousser plus
loin mes propres recherches.

Après avoir fait choix d'un des points de la surface, on peut passer de ce
point à un point voisin dans une infinité de directions. Il y a bien là
autant de courbures différentes ; mais, quelque soit la diversité des
configurations, la répartition de la courbure autour du point choisi offre un
arrangement symétrique.

Les propriétés des courbures principales ont été examinées avec beaucoup de
soin ; celles des courbures moyennes n'ont même pas été remarquées. Le but de
ce mémoire sera de les faire connoître et d'en déduire la loi suivant
laquelle la courbure est répartie autour d'un des points de la surface.

Le lecteur verra sans doute avec quelque surprise une loi qu'il pourroit être
tenté de regarder comme illusoire être, pour ainsi dire, calquée sur une
autre loi qui doit lui être déjà familière, celle de la distribution de la
chaleur dans l'anneau.

Cette analogie inattendue n'est cependant pas purement fortuite : on s'en
rendra raison en observant que, de part et d'autre, l'emploi des quantités
moyennes se trouve uni à la considération d'un espace circulaire.

Un préjugé favorable à la notion des moyennes courbures paroît devoir naître
d'un tel rapprochement ;

\page{46}
pourroit-on penser, en effet, que des remarques fantastiques vinssent se
ranger, d'elles mêmes, dans un ordre que l'on sait appartenir ailleurs à des
vérités démontrées ?

Je suivrai dans tous ses détails la comparaison qui vient d'être indiquée.
Elle offrira aux géomètres l'occasion de se familiariser avec un genre de
quantités demeuré jusqu'à présent étranger à leurs recherches ; et,
peut-être, à la faveur d'une marche en quelque sorte parallèle,
obtiendrai-je pour mes propres travaux une foible partie de l'attention qui
ne fut jamais refusée à l'auteur dont j'aurai tenté de suivre les traces.

La notion des courbures moyennes dérive de la comparaison entre la courbure
de la surface et celle de la sphère. Une construction imaginée dans la vue de
rendre cette comparaison sensible conduit aux propositions suivantes :

La courbure de la sphère est en raison directe de la surface convexe d'un
petit cylindre à base circulaire compris entre le plan tangent
\struck{au point} choisi et celui qui couperoit la sphère suivant le cercle
qui sert de base au cylindre :

Lorsque les rayons de courbures principales sont dirigés du même côté, la
courbure de la surface est en raison directe de la surface convexe d'un
segment cylindrique à base d'ellipse, appuyé sur le plan tangent au point
choisi, et terminé, à l'autre extrémité, par une courbe rentrante à double
courbure, composée de quatre parties symétriques et tracée, sur la surface,
autour du point de tangence.

La surface convexe du segment cylindrique et celle

\page{47}
du petit cylindre qui appartient à la sphère de \emph{courbure moyenne} sont
d'une étendue égale : et, par conséquent, la courbure de la surface est en
raison directe de la surface convexe de ce cylindre.

On verra dans la suite qu'une proposition analogue, mais un peu plus
compliquée, est applicable au cas où les deux rayons de courbures principales
de la surface ne sont pas dirigés du même côté.

Suivant cette manière d'envisager les choses, la figure et la quantité de
courbure qui appartiennent à la surface obtiennent des représentations
séparées ; car, tandis que la première détermine la figure du segment
cylindrique, la seconde, a pour mesure l'étendue de la superficie courbée du
même segment, étendue qui, ainsi qu'on vient de le dire, est égale à celle de
la petite surface cylindrique donnée par la sphère de \emph{courbure
moyenne}.

À l'aide de ces remarques, il sera facile de concevoir comment, dans les
questions où il s'agit d'apprécier la force qui dépend de la courbure d'une
surface, on trouve que cette force ne dépend que de la courbure moyenne.

Si on admet, en effet, que chacun des points de la surface du petit cylindre
représente des parties égales de cette même force, on sera mené à dire que la
force totale est proportionnelle à l'étendue de la surface entière de ce
petit cylindre.

Ainsi la notion fondamentale des courbures moyennes, à laquelle on avoit
d'abord été conduit par l'examen de questions empruntées à la mécanique,
fournira, revêtue ensuite des formes géométriques qui lui conviennent,
l'explication des relations dont le premier aperçu a rendu l'existence de ce
genre de quantités

\note{les deux dernières lignes de la page sont écrites très petit, au ras du
bord inférieur, et la phrase s'arrête là, sans son dernier mot ; un double
trait horizontal la suit. La page 6 recommence par un titre, et le mémoire ne
reprend pas cette phrase.}

\section{Mémoire sur la courbure des surfaces}

\page{48}
N° 1 Si, par rapport aux surfaces, on avoit besoin de connoître la mesure de
la courbure, on trouveroit peu de secours dans les écrits des géomètres qui
se sont occupés des diverses questions dont se compose la géométrie
descriptive ; et l'on seroit forcé de recourir à des travaux plus anciens.

Le mémoire d'Euler intitulé : \emph{Recherches sur la courbure des surfaces}
(M. de Berlin 1760 p. 119) contient, en effet, tout ce que l'on sait
d'important à cet égard ; et, en lisant avec attention ce beau mémoire, on
apperçoit bientôt que l'illustre auteur y a déposé le germe des recherches
qui peuvent faire disparoître les difficultés qu'il a pris soin de signaler.

Dans la vue de fixer le point de départ on reproduira ici les premiers
passages de ce mémoire.

En voici le début :

\begin{quote}
« Pour connoître la courbure des lignes courbes, la détermination du rayon
osculateur en fournit la plus juste mesure, en nous présentant pour chaque
point de la courbe un cercle, dont la courbure est précisément la même. Mais,
quand on demande la courbure d'une surface, la question est fort équivoque,
et point du tout susceptible d'une réponse absolue, comme dans le cas
précédent. Il n'y a que les surfaces sphériques dont on puisse mesurer la
courbure, attendu que la courbure d'une sphère est la même que celle de ses
grands cercles et que son rayon en peut être regardé comme la juste mesure.
Mais pour les autres surfaces, on n'en sauroit même comparer la courbure avec
celle
\end{quote}

\page{49}
\begin{quote}
d'une sphère, comme on peut toujours comparer la courbure d'une ligne courbe
avec celle d'un cercle ; la raison en est évidente puisque, dans chaque point
d'une surface, il peut y avoir une infinité de courbures différentes\ldots\ »
\end{quote}

Plus bas (p. 120) on lit le passage suivant

\begin{quote}
« Mais, pour le sujet présent, il suffit de ne considérer de toutes ces
infinies sections que celles qui sont perpendiculaires à la surface, dont le
nombre est pourtant encore infini\ldots, et l'assemblage de tous ces rayons
(ceux des courbes produites par les sections normales) nous donnera la juste
mesure de la courbure de la surface au point donné ».
\end{quote}

Lorsque je proposai de représenter la courbure d'une surface par la somme des
raisons inverses des deux rayons de ses courbures principales, l'objection
indiquée d'avance dans le premier des passages qu'on vient de lire fut
reproduite ; et les opposans parurent croire que, parmi le nombre infini des
courbes contenues dans les différens plans menés par un des points de la
surface, le choix des courbures principales étoit arbitraire.

Dans le mémoire envoyé à l'académie, pour le concours de 1816, je pris la
somme des raisons inverses des rayons de courbures qui résultent de toutes
les intersections que produisent les différentes positions du plan mené par
le point donné.
\struck{Or, soit que la formule qui exprime la courbure, elle}
\struck{Cette somme se réduit d'abord à celle qui résulteroit des seules
sections normales, et cette dernière donne la somme} des raisons inverses des
rayons des courbures principales, répétée un nombre infini de fois ; je
conclus enfin que la somme totale n'étoit autre chose que le résultat
d'opérations équivalentes dont une seule exprimeroit la mesure de la
courbure ; opération qui, suivant cette manière d'envisager les choses, a dû
être répétée un nombre infini de fois.

\marginal{ces sections obliques, renferme le cosinus de l'inclinaison. En
prenant par rapport à chacune des positions du plan normal la somme de toutes
les sections obliques}

\note{la phrase biffée porte deux couches : « Cette somme se réduit\ldots » est
la première, et au dessus, dans l'interligne, une reprise qui commence par un
mot lu ici « Or, soit » sans certitude ; l'une et l'autre sont barrées. Un
signe « + » est porté au dessus de « la somme » ; le même signe reparaît sur
le béquet de la vue 52, dont le texte pourrait être le remplacement prévu
pour ce passage — la vue seule ne permet pas de l'affirmer. La note marginale
est écrite verticalement dans la marge de gauche, en regard de ces lignes.}

\page{50}
Un pareil raisonnement étoit conforme à l'indication donnée par Euler ; car
on vient de voir que cet habile géomètre pensoit, sans qu'il ait pourtant
témoigné l'avoir vérifiée, qu'on pouvoit mettre à l'écart les sections
obliques et que la somme des sections normales conduiroit à connoître « la
juste mesure de la courbure de la surface au point donné ». Il étoit
d'ailleurs évident que l'idée de l'infini, étrangère à une semblable mesure,
y avoit été introduite par la multiplicité des opérations ; multiplicité à
laquelle avoit donné lieu la forme du raisonnement, et non la nature de la
question.

Quoiqu'il en soit du mérite de cette démonstration, elle fut formellement
désapprouvée ; et, la légitimité de l'hypothèse ne cessant pas d'être
contestée, il parut nécessaire d'examiner de nouveau quelle est, par rapport
aux surfaces, l'expression de la courbure.

En faisant usage des formules d'Euler, il est facile de voir que, si on
substitue, aux plans qui contiennent les principales courbures de la surface,
deux autres plans normaux, perpendiculaires entr'eux, la somme des raisons
inverses des rayons de courbures des courbes contenues dans ces deux plans,
est égale à celle des raisons inverses des \ill{} rayons de courbures
principales.

Cette observation devoit se présenter la première ; aussi fut-elle employée
dans la démonstration qui vient d'être rappelée.

De nouvelles réflexions sur le même sujet firent naître les remarques
suivantes :

Si, parmi les systèmes de deux plans normaux perpendiculaires entr'eux, on
choisit ceux qui font d'angle de $45^{\circ}$ avec les plans qui contiennent
les

\note{une tache d'encre couvre entièrement un mot avant « rayons de courbures
principales » ; elle n'a pas été devinée. Un béquet dépasse au bas de la vue,
dont l'écriture transparaît à l'envers : c'est celui que la vue 52 montre de
face.}

\page{51}
courbures principales, un même rayon appartient aux courbes formées par les
intersections de chacun des plans choisis.

Les quatre quadrans de la sphère décrite de ce rayon sont situés
alternativement, au dehors et au dedans de la surface :

Enfin \ill{}, par une suite de la manière dont la courbure est distribuée
autour d'un point pris sur la surface, il s'établit une compensation parfaite
entre les courbures plus grandes et moindres que celle de la sphère ; ensorte
que, si on prend, dans un ordre facile à déterminer, un nombre donné des
raisons inverses qui mesurent ces courbures, leur somme sera égale à la
raison inverse du rayon de courbure de cette sphère, multipliée par le nombre
donné.

Les conséquences principales de ces observations suffisent certainement pour
établir la notion des courbures moyennes, et elles ont été déjà publiées,
néanmoins, les explications données jusqu'à présent, mêlées, \ill{}, avec
d'autres recherches qui peuvent détourner l'attention, sont \ill{}

Pour mettre le lecteur à même de se former une idée complète de ce genre de
quantités, il convient de le montrer que la question des courbures moyennes
est indépendante des applications qu'on en peut faire. On se propose ici de
prendre cette question pour l'objet d'un travail spécial, et de lui donner
tous les développemens dont elle est susceptible.

N° 2 Quand on remonte aux formules d'Euler, on voit que, si $f$ et $g$ sont
le plus grand et le moindre rayons

\note{« Enfin » est suivi d'un chiffre abrégé — 3 ou 5 — que la vue ne permet
pas de trancher, et aucune énumération antérieure n'est visible sur la page ;
il n'a pas été deviné. « distribuée autour » est ajouté dans la marge de
gauche, à la hauteur de la ligne où il s'insère. Deux amas de mots biffés et
récrits l'un sur l'autre, l'un après « mêlées », l'autre à la fin de la même
phrase, n'ont pas été lus.}

\page{52}
chaque valeur du cosinus est écrite avec des signes opposés, cette somme
disparoît donc et, par conséquent il ne reste plus, dans la somme générale,
que les seules valeurs dues aux sections normales. Il est ensuite facile de
voir que cette dernière somme est composée de celle \struck{des raisons
inverses}

\note{cette vue ne montre qu'un béquet, posé seul sur un fond blanc : une
bande de papier portant six lignes, dont la dernière est biffée, et un signe
« + » au dessus de « celle ». Le même béquet dépasse, vu de dos et inversé,
au bas de la vue 50. Le signe « + » et le sujet du passage — la disparition
des sections obliques dans la somme générale — le rapprochent du passage
biffé de la vue 49 ; la vue ne dit pas où il devoit être collé.}

\page{53}
\note{seconde photographie de la page transcrite à la vue 51 : même feuillet,
même numéro à l'encre 26, même pagination 9, mêmes ratures et même ajout
marginal. Elle n'est pas transcrite une seconde fois.}

\page{54}
osculateurs, et $r$ le rayon osculateur qui appartient à la courbe due à la
section \ill{} plan normal fesant l'angle $\varphi$
avec celui de plus grande courbure, on aura
\[
  r = \frac{2fg}{f + g - (f-g)\cos 2\varphi} .
\]
De même aussi, $r'$ étant le rayon osculateur dû à la section par le plan
normal qui fait l'angle $\varphi + \frac{\pi}{2}$ avec celui de plus grande
courbure, on aura
\[
  r' = \frac{2fg}{f + g - (f-g)\cos (2\varphi + \pi)}
     = \frac{2fg}{f + g + (f-g)\cos 2\varphi} .
\]
Ces formules donnent :
\begin{align*}
  \frac{1}{r} &= \frac{f + g - (f-g)\cos 2\varphi}{2fg}\\
     &= \frac{(f+g)\{\cos^{2}\varphi + \sin^{2}\varphi\}
               - (f-g)(\cos^{2}\varphi - \sin^{2}\varphi)}{2fg}\\
     &= \frac{f\sin^{2}\varphi + g\cos^{2}\varphi}{fg} ;
\end{align*}
\begin{align*}
  \frac{1}{r'} &= \frac{f + g + (f-g)\cos 2\varphi}{2fg}
     = \frac{f\cos^{2}\varphi + g\sin^{2}\varphi}{fg} .
\end{align*}
En considérant les diverses sections normales d'où naissent autant de
courbures différentes, la forme de la valeur des raisons $\frac{1}{r}$ et
$\frac{1}{r'}$, dues à deux sections normales perpendiculaires entr'elles,
conduit aisément à reconnoître quelle est la loi de la distribution de la
courbure autour d'un point donné de la surface ; mais, avant d'aller plus
loin, il convient de prévenir une objection qui se présente assez
naturellement.

Il s'agit de la courbure des surfaces, et pourtant, dans les recherches qui
vont suivre immédiatement, les seules formules employées appartiennent à la
courbure linéaire. À la vérité, on sent que ces deux genres de courbures ont
entr'elles une liaison nécessaire ; mais la nature de cette liaison demeure
indéterminée.

\note{un ou deux mots biffés et récrits l'un sur l'autre précèdent « plan
normal » à la deuxième ligne ; ils n'ont pas été lus. Les deux dernières
chaînes d'égalités sont disposées côte à côte sur le feuillet, celle de $1/r$
à gauche et celle de $1/r'$ à droite ; elles sont mises l'une sous l'autre
ici. Dans la première, le numérateur $2fg$ a été écrit seul puis biffé avant
que la fraction ne soit récrite ; une tache d'encre couvre le chiffre 1 de
$1/r$ sans le rendre douteux.}

\page{55}
La difficulté qui vient d'être signalée ne doit cependant pas arrêter le
lecteur : elle résulte uniquement de la manière dont on envisage les
questions ; et peut-être seroit-il embarrassant de l'éviter, lorsqu'on se
propose d'exclure, comme on a crû devoir le faire, d'abord, toute
considération inusitée.

Dans la suite de ce mémoire, lorsque moins de réserve paroîtra nécessaire, il
sera démontré qu'il existe la plus grande analogie entre la manière dont on
mesure la courbure linéaire et celle qui peut servir à mesurer la courbure
moyenne des surfaces.

N° 3 La loi suivant laquelle la courbure est distribuée entre les différentes
intersections produites par les plans normaux menés par un des points de la
surface, est renfermée dans un petit nombre de propositions ; et, comme on
l'a déjà dit ces propositions, comparées à celles qui s'appliquent à la
distribution de la chaleur dans l'anneau, offrent une analogie remarquable.
Mais on saisiroit mal les ressemblances annoncées si, avant tout, nous ne
fesions remarquer que les deux lois présentent cependant une entière
différence, par rapport à l'intervalle qui sépare les points où existent les
propriétés extrêmes ; ainsi, d'une part, lorsqu'il s'agit de la courbure, la
distance angulaire entre la plus grande et la moindre, contenues dans des
plans perpendiculaires entr'eux, a pour mesure un quart de circonférence :
tandis que, de l'autre, par rapport à l'anneau, la plus et la moins élevée
des températures appartiennent à deux points séparés par une demi
circonférence. Il est facile de prévoir les conséquences de

\page{56}
cette différence ; nous aurons soin d'en faire mention, et on pourra se
convaincre qu'elles ne dérangent en rien la comparaison que nous avons en
vue.

Nous commencerons par énoncer et démontrer sommairement chacune des
propositions dont se compose la loi de la distribution de la courbure.
Mettant ensuite sous les yeux du lecteur la proposition correspondante
empruntée à la théorie de la chaleur, nous chercherons dans l'imitation d'un
modèle si parfait, en même tems qu'une nouvelle expression de notre
proposition, la preuve certaine de l'analogie qui doit servir de
recommandation à nos propres recherches.

1\textsuperscript{re} prop. \emph{Quelle que soit d'ailleurs la position de
deux plans normaux perpendiculaires entr'eux menés par un point donné de la
surface, la somme des courbures contenues dans ces deux plans sera toujours
la même : par conséquent, cette somme sera égale à celle des courbures
principales.}

Le raisonnement suivant \struck{suffit} pour établir la vérité de cette
proposition :

La courbure linéaire est exprimée par la raison inverse du rayon osculateur ;
$r$ et $r'$ étant les rayons osculateurs des courbes contenues dans deux
plans normaux perpendiculaires entr'eux, les formules
$\frac{1}{r} = \frac{f\sin^{2}\varphi + g\cos^{2}\varphi}{fg}$,
$\frac{1}{r'} = \frac{f\cos^{2}\varphi + g\sin^{2}\varphi}{fg}$, données
N° 2, conduisent à l'équation
$\frac{1}{r} + \frac{1}{r'} = \frac{1}{f} + \frac{1}{g}$.

Cette équation est indépendante de la valeur de l'angle $\varphi$ ; elle est
donc indépendante de la position du système des plans choisis.

En n'écrivant ni les formules ni les développemens qui s'y rapportent, voici,
à l'égard de l'anneau, la proposition correspondante : (V. Théorie de la
chaleur par Mr Fourier \ill{} IV N° 247 p. 277-27\ill{}

\note{la dernière ligne de la page court jusqu'au bord du feuillet et s'y
perd : ni l'abréviation qui précède « IV », ni les derniers chiffres de la
pagination n'ont été lus. La vue 58 cite le même livre « (p. 278-279) ».}

\page{57}
\begin{quote}
« Si, dans cet état, (celui qui s'établit au bout d'un certain tems) on
choisit deux points de l'anneau, situés aux deux extrémités d'un même
diamètre\ldots\ la demi somme des températures des points opposés sera une
quantité\ldots\ qui seroit encore la même si on avoit choisi deux points
situés aux extrémités d'un autre diamètre ».
\end{quote}

Nous pouvons dire aussi : si parmi les courbes qui passent par le point
donné, on en choisit deux comprises, l'une dans un \struck{premier} plan
normal, l'autre dans un second plan normal, perpendiculaire au premier, la
demi somme des courbures sera une quantité $\frac{1}{f} + \frac{1}{g}$, qui
seroit encore la même si on avoit choisi les courbes comprises dans deux
\struck{autres} plans normaux perpendiculaires entr'eux.

Ainsi, de part et d'autre, l'intervalle qui sépare les quantités extrêmes,
est le même que celui qui sépare des quantités différentes, quand celles-ci
satisfont à cette condition : que leur somme soit égale à celle des mêmes
quantités extrêmes.

2\textsuperscript{ème} prop. \emph{La courbure contenue dans le plan normal
qui fait l'angle de $45^{\circ}$ avec ceux qui contiennent la plus grande et
la moindre courbure est moyenne entre les courbures principales de la
surface.}

Cette proposition résulte des formules dont nous avons déjà fait usage. En
effet, $r$ étant le rayon osculateur de la courbe contenue dans le plan
normal qui fait l'angle $\varphi$ avec celui de plus grande courbure, on a
$\frac{1}{r} = \frac{f\sin^{2}\varphi + g\cos^{2}\varphi}{fg}$ : quand
$\varphi = 45^{\circ}$, cette formule donne
$\frac{1}{r} = \frac{f+g}{2fg} = \frac{1}{2}\left(\frac{1}{f} +
\frac{1}{g}\right)$ : la courbure exprimée par $\frac{1}{r}$ est donc alors
moyenne entre les courbures principales, représentées par $\frac{1}{f}$ et
$\frac{1}{g}$.

\note{la quantité appelée « demi somme » est écrite $\frac{1}{f} +
\frac{1}{g}$ sur le feuillet, sans le facteur $\frac{1}{2}$ que la suite de
la page emploie ; la lecture est sûre et n'a pas été corrigée.}

\page{58}
Il est, au reste, évident que, dans le cas dont nous parlons, les courbes
contenues dans les deux plans normaux perpendiculaires entr'eux sont
semblables ; et, s'il étoit besoin de le démontrer, on diroit que la formule
$\frac{1}{r'} = \frac{f\cos^{2}\varphi + g\sin^{2}\varphi}{fg} =
\frac{f+g}{2fg}$, relative à la position du second des plans normaux dont il
s'agit, conduiroit en effet aussi à conclure que la valeur de $\frac{1}{r'}$
est moyenne entre celles de $\frac{1}{f}$ et de $\frac{1}{g}$.

En rapprochant les passages que séparent des développemens propres à
l'ouvrage, on peut extraire du livre déjà cité (p. 278-279), le texte
suivant, qui, par rapport à la distribution de la chaleur dans l'anneau,
renferme les conséquences de la proposition analogue à celle qui nous occupe.

\begin{quote}
« Si on cherche d'abord le point de l'anneau pour lequel on a la
condition\ldots\ (condition qui est remplie à l'égard du point qui sépare en
deux parties égales l'arc de 180 compris entre les points auxquels
appartiennent les températures extrêmes), on voit que la température de ce
point est à chaque instant la température moyenne de l'anneau : il en est de
même du point diamétralement opposé :\ldots\ ces deux points divisent la
circonférence de l'anneau en deux parties dont l'état est pareil, mais de
signe opposé : chaque point de l'une de ces parties a une température qui
excède la température moyenne, et \struck{la quantité} de cet excès est
proportionnelle au sinus\ldots\ chaque point de l'autre partie a une
température moindre que la température moyenne et la différence est la même
que l'excès dans le point opposé ».
\end{quote}

\note{la parenthèse « (condition qui est remplie\ldots) » est de Germain et
non de Fourier : elle est écrite sans guillemet, au milieu de la citation.}

\page{59}
Par rapport à la courbure :

Si on cherche d'abord la position du plan qui donne lieu à la condition
$\frac{1}{r} = \frac{f+g}{2fg}$ (condition qui est remplie à l'égard de la
courbe contenue dans le plan qui sépare en deux parties égales l'angle de 90,
compris entre les plans des courbures extrêmes ou principales de la surface,)
on démontrera que la courbure contenue dans ce plan est égale à la courbure
moyenne de la surface : il en est de même de la courbure contenue dans le
plan normal perpendiculaire au premier. Ces deux plans, prolongés en deça et
au delà du point donné, divisent le contour de ce point en quatre
\uncertain{régions} dont l'état est pareil mais alternativement de signe
opposé. Chaque courbe située dans l'une de ces parties a une courbure qui
excède la courbure moyenne, et la quantité de cet excès est proportionnelle à
$\frac{1}{r} - \frac{f+g}{2fg}$.

Chaque courbe située dans une des \uncertain{régions} contiguës à la première
a une courbure moindre que la courbure moyenne et la différence est la même
que l'excès dans la partie voisine.

L'analogie se fait remarquer ici en ce que le lieu, soit du point, soit de la
courbe, qui est en possession de la valeur moyenne, occupe le milieu de la
distance angulaire qui sépare ou les points, ou les courbes, auxquelles
appartiennent les valeurs extrêmes. L'analogie est encore \struck{sensible}
sensible relativement à l'état des points ou des courbes répartis dans les
deux, ou dans les quatre parties que l'on compare ; car, de part et d'autre,
cet état offre des rapports semblables. On voit aussi que le nombre de ces
parties dépend, dans les lois que nous \ill{} de l'intervalle qui sépare les
valeurs extrêmes.

\note{le mot lu « régions », deux fois, est écrit d'un trait rapide et
repassé ; la lecture reste douteuse. À l'avant-dernière ligne, deux ou trois
mots biffés et couverts d'encre après « dans les lois que nous » n'ont pas
été lus.}

\page{60}
Pour connoître la répartition de la courbure autour d'un des points de la
surface, il suffit sans doute d'examiner la moitié du contour de ce point. On
trouve cependant quelqu'avantage à considérer, comme on l'a fait dans ce qui
précède, le contour entier : alors, le point donné étant l'origine des
différentes courbes, nous dirons qu'un plan normal prolongé en deça et au
delà du point donné contient deux courbes semblables.

Cela posé, voici la proposition dont il s'agit présentement :

3\textsuperscript{ème} prop. \emph{Il existe, en général, autour du point
donné, quatre directions suivant lesquelles les courbes contenues dans les
plans normaux sont semblables. À l'égard des courbures principales, chacune
est contenue dans un seul plan, et, par conséquent, le nombre des courbes
semblables se réduit à deux.}

Une remarque fort simple suffit à la démonstration de cette proposition :

En vertu des formules connues, on a
\begin{align*}
  \sin (45^{\circ} + \varphi') &= \tfrac{1}{2}(\cos.\varphi' + \sin.\varphi')\\
  \sin \left(\tfrac{\pi}{2} + 45^{\circ} - \varphi'\right)
    &= \tfrac{1}{2}(\cos.\varphi' + \sin.\varphi')\\
  \cos (45^{\circ} + \varphi') &= \tfrac{1}{2}(\cos.\varphi' - \sin.\varphi')\\
  \cos \left(\tfrac{\pi}{2} + 45^{\circ} - \varphi'\right)
    &= \tfrac{1}{2}(\cos.\varphi' - \sin.\varphi')
\end{align*}
Deux valeurs différentes de $\varphi$, savoir $\varphi = 45^{\circ} +
\varphi'$, $\varphi = \frac{\pi}{2} + 45^{\circ} - \varphi'$, donnent donc
également
\begin{align*}
  \frac{1}{r} &= \frac{f\sin^{2}\varphi + g\cos^{2}\varphi}{fg}\\
    &= \frac{f(\cos.\varphi' + \sin.\varphi')^{2}
             + g(\cos.\varphi' - \sin.\varphi')^{2}}{2fg}\\
    &= \frac{f + g + (f-g)\sin. 2\varphi'}{2fg} .
\end{align*}
En prenant $\varphi' = 45$, les deux valeurs de $\varphi$ se réduisent à une
seule, et le second membre de l'équation
$\frac{1}{r} = \frac{f + g + (f-g)\sin. 2\varphi'}{2fg}$ devient
$\frac{1}{f}$. Ainsi, à l'égard de la détermination des courbes

\note{les quatre identités sont écrites avec le facteur $\frac{1}{2}$ ; le
calcul qui les suit immédiatement les emploie comme si ce facteur étoit
$\frac{\sqrt{2}}{2}$. Aucun radical n'est visible sur la vue, et rien n'a été
corrigé ici. La vue 60 s'arrête au milieu de la phrase ; le texte se poursuit
à la vue 61, hors de ce lot.}

\end{document}
